\section*{第 \ref{ch:b1:diagenesis-basins} 章——沉积成岩作用与沉积盆地}

\begin{solution}{exo:b1:diagenesis-basins:1}
砂岩 $0.49\,e^{-2000/3704} = 0.49\times0.583 = \num{0.286}$；页岩
$0.63\,e^{-2000/1960} = 0.63\times0.360 = \num{0.227}$。孔隙度减半时
$e^{-z/\lambda} = 1/2$，即 $z = \lambda\ln 2$：砂岩为 \qty{2567}{m}，
页岩为 \qty{1359}{m}。
\end{solution}

\begin{solution}{exo:b1:diagenesis-basins:2}
$\Delta A = \Delta E + S = 30 + 50 = \qty{80}{m}$。其中 \qty{60}{m} 由
沉积物充填，故 $\Delta W = 80 - 60 = \qty{20}{m}$：陆架加深
\qty{20}{m}。
\end{solution}

\begin{solution}{exo:b1:diagenesis-basins:3}
$h = Y(\rho_m - \rho_w)/(\rho_m - \rho_s)$：取 \qty{2300}{kg.m^{-3}} 时，
$1.0\times2270/1000 = \qty{2.27}{km}$；取 \qty{2500}{kg.m^{-3}} 时，
$1.0\times2270/800 = \qty{2.84}{km}$。每一米沉积物都取代柱体中一米的\omterm{def:b1:layered-earth:layers}{地幔}，
沉积物越致密，与\omterm{def:b1:layered-earth:layers}{地幔}相比所缺的质量越小，
因此要补足与充水空洞相同的亏损，所需的
沉积物就越多。
\end{solution}

\begin{solution}{exo:b1:diagenesis-basins:4}
(a)~方解石 $100.09/2.71 = \qty{36.93}{cm^3.mol^{-1}}$；白云石 $184.40/2.84
= \qty{64.93}{cm^3.mol^{-1}}$。(b)~两摩尔方解石占据
\qty{73.86}{cm^3}，一摩尔白云石占 \qty{64.93}{cm^3}，固体体积的收缩比例为
$1 - 64.93/73.86 = \qty{12.1}{\%}$。(c)~在一立方米中，\qty{0.90}{m^3}
的方解石变为 $0.90\times(1 - 0.121) = \qty{0.791}{m^3}$ 的白云石；
孔隙度变为 $1 - 0.791 = \num{0.21}$。
\end{solution}

\begin{solution}{exo:b1:diagenesis-basins:5}
(a)~$s = 1000 - 0.63\times1960\times(1 - e^{-1000/1960}) = \qty{506.5}{m}$；
平均孔隙度为 $1 - 506.5/1000 = \num{0.493}$。(b)~顶面位于
\qty{2000}{m} 处，由 $s(2000, b) = 506.5$ 解得 $b = \qty{2628.7}{m}$：该层
厚 \qty{628.7}{m}，平均孔隙度为 $1 - 506.5/628.7 = \num{0.194}$。
(c)~颗粒不变，因此失去的水量等于失去的厚度：
$1000 - 628.7 = \qty{371}{m^3}$，即每平方米损失的水量。
\end{solution}

\begin{solution}{exo:b1:diagenesis-basins:6}
(a)~$s = 300 - 0.49\times3704\times(e^{-2500/3704} - e^{-2800/3704}) =
\qty{228.10}{m}$。(b)~由 $b_0 = 300$ 出发：$s(0, 300) = 158.80$，$1 - \varphi(300) =
0.5481$，$b_1 = 300 + (228.10 - 158.80)/0.5481 = \qty{426.44}{m}$；然后
$s(0, 426.44) = 229.07$，$1 - \varphi = 0.5633$，$b_2 = 426.44 - 0.96/0.5633 =
\qty{424.73}{m}$；第三步的改变量不足一毫米。砂岩
在地表时厚 \qty{424.7}{m}。(c)~$424.7/300 = 1.42$：埋藏
使其减薄了 \num{1.42} 倍。
\end{solution}

\begin{solution}{exo:b1:diagenesis-basins:7}
(a)~冰中所含水量为 $52\times10^6\times916.7/1000 = \qty{4.77e7}{km^3}$；
将其平铺于 \qty{361.84e6}{km^2} 的面积上，折合 $\qty{0.1317}{km} = \qty{132}{m}$ 的海平面
上升。(b)~由海岸线得到的值为 \qty{134}{m}，多 \qty{2}{m}。(c)~$4.767\times
10^7/0.134 = \qty{355.7e6}{km^2}$，比今天少 \qty{1.7}{\%}。海平面
下降时大陆架出露，因此冰期海洋面积较小，而
从中取走同样体积的水，会使海平面下降得更多。
\end{solution}

\begin{solution}{exo:b1:diagenesis-basins:8}
(a)~$(-9.1 - (-98.1))/8 = \qty{11.1}{m}$ 每千年，接近
由海岸线得出的 \qty{12}{m}。(b)~每年 \qty{40}{mm}，即每千年 \qty{40}{m}，
是平均值的 $40/12 = 3.3$ 倍；五百年内，
变化量为 $0.040\times500 = \qty{20}{m}$。
\end{solution}

\begin{solution}{exo:b1:diagenesis-basins:9}
(a)~$S = 50\times2 = \qty{100}{m}$，$\Delta E = \qty{-40}{m}$：
$\Delta A = \qty{60}{m}$，因水深未变，全部由沉积物充填。
(b)~在 \num{20000} 年内，基底仅\omterm{def:b1:diagenesis-basins:subsidence}{沉降} $50\times0.02 = \qty{1}{m}$：
\omterm{def:b1:diagenesis-basins:eustasy}{相对海平面}下降 $120 - 1 = \qty{119}{m}$。海底原位于
海平面以下 \qty{20}{m}，最终高出海平面 $119 - 20 = \qty{99}{m}$：整个
陆架全部出露。
\end{solution}

\begin{solution}{exo:b1:diagenesis-basins:10}
(a)~对于薄层，孔隙度在层内近似不变：其固体厚度
在深度处为 $s \approx h(1 - \varphi(z))$，在地表处为 $s \approx h_0(1 -
\varphi_0)$。令二者相等，得 $h_0 \approx h(1 -
\varphi(z))/(1 - \varphi_0)$。(b)~$\varphi(3000) = 0.63\,e^{-3000/1960} =
0.136$，于是 $h_0 \approx 10\times0.864/0.37 = \qty{23.3}{m}$，与
精确值 \qty{23.1}{m} 相当接近。(c)~恢复到地表后，该层
厚 \qty{23}{m}，相对于地表附近孔隙度的衰减已不能视为薄层：
其孔隙度由顶部的 $\varphi_0$ 降至底部更小，因此
其颗粒所占厚度略小于
按 $\varphi_0$ 全层取值算得的厚度。
\end{solution}

\begin{solution}{exo:b1:diagenesis-basins:11}
(a)~$s = 3000 - 0.63\times1960\times(1 - e^{-3000/1960}) = \qty{2032}{m}$；平均孔隙度为
$1 - 2032/3000 = \num{0.323}$，平均密度为
$0.323\times1030 + 0.677\times2700 = \qty{2161}{kg.m^{-3}}$。(b)~因子为
$(3300 - 1030)/(3300 - 2161) = \num{1.99}$；充填物相当于受水载荷的
\omterm{def:b1:diagenesis-basins:subsidence}{沉降}，\omterm{def:b1:diagenesis-basins:subsidence}{沉降}量为 $3000/1.99 = \qty{1505}{m}$。(c)~若采用颗粒密度，则因子为
$2270/600 = \num{3.78}$，\omterm{def:b1:diagenesis-basins:subsidence}{沉降}量仅为 \qty{793}{m}。这样做是
错误的，因为三分之一的充填物是水：沉积物远轻于其
颗粒的密度，质量平衡中计入的只有
饱和沉积物的密度。
\end{solution}

\begin{solution}{exo:b1:diagenesis-basins:12}
(a)~该立方体中含方解石 \qty{0.90}{m^3} $= \qty{9.0e5}{cm^3}$，即
$9.0\times10^5/36.93 = \num{24370}$ 摩尔。(b)~每两摩尔方解石对应一摩尔镁：
\num{12180} 摩尔。(c)~每千克海水含镁 $35.165\times0.0365 =
\qty{1.284}{g}$，即 $1.284/24.305 = \qty{0.0528}{mol}$；反应
需要 $12180/0.0528 = \qty{2.31e5}{kg}$ 海水，即 $\qty{224}{m^3}$——
超过岩石体积的两百多倍，即使镁离子全部
被利用。(d)~白云石大规模形成，只可能发生在极大量的
水长时间流经岩石的地方，或水为
镁含量高于海水的卤水：例如在被下沉卤水冲刷的
台地中，而非在保持同一孔隙水的沉积物中。
\end{solution}

\begin{solution}{pb:b1:diagenesis-basins:1}
\textbf{1.} 砂岩：$\varphi(0) = \num{0.490}$，$\varphi(800) =
0.49\,e^{-800/3704} = \num{0.395}$。页岩：$\varphi(800) = \num{0.419}$，
$\varphi(2600) = \num{0.167}$。白垩：$\varphi(2600) = \num{0.110}$，
$\varphi(3200) = \num{0.072}$。

\textbf{2.} 取自\cref{prop:b1:diagenesis-basins:solid}的 $s$，且 $\bar\varphi
= 1 - s/h$（问题~7）：砂岩 \num{0.441}，页岩 \num{0.274}，白垩
\num{0.090}。

\textbf{3.} $h\bar\varphi$：砂岩 \qty{353}{m^3}，页岩 \qty{493}{m^3}，白垩
\qty{54}{m^3}，均为每平方米的水量。

\textbf{4.} $\bar\varphi\rho_w + (1 - \bar\varphi)\rho_g$：砂岩
$0.441\times1030 + 0.559\times2650 = \qty{1936}{kg.m^{-3}}$，页岩
\qty{2242}{kg.m^{-3}}，白垩 \qty{2559}{kg.m^{-3}}。该柱体的质量为
$1936\times800 + 2242\times1800 + 2559\times600 = \qty{7.12e6}{kg}$，即每平方米
质量，平均密度为 \qty{2225}{kg.m^{-3}}。

\textbf{5.} \omterm{def:b1:stress-strain-rheology:lithostatic}{静岩压力} $2225\times9.81\times3200 = \qty{69.9}{MPa}$；静水压力
$1030\times9.81\times3200 = \qty{32.3}{MPa}$；\omterm{def:b1:stress-strain-rheology:lithostatic}{有效应力}
$69.9 - 32.3 = \qty{37.5}{MPa}$。

\textbf{6.} 由 $0.49\,e^{-z/3704} = 0.63\,e^{-z/1960}$ 得 $z = \ln(0.63/0.49)
/(1/1960 - 1/3704) = \qty{1046}{m}$。在此深度以下砂岩的孔隙度更大：
页岩起始含水更多，但失水所经历的距离更短。

\textbf{7.} 砂岩 $800 - 0.49\times3704\times(1 - e^{-800/3704}) =
\qty{447.4}{m}$；页岩 \qty{1306.7}{m}；白垩 \qty{546.0}{m}。

\textbf{8.} 颗粒部分厚 \qty{2300.2}{m}；水占柱体的 $1 - 2300.2/3200 =
\qty{28}{\%}$。

\textbf{9.} 砂岩 $447.4/15 = \qty{29.8}{m}$，页岩 $1306.7/40 =
\qty{32.7}{m}$，白垩 $546.0/20 = \qty{27.3}{m}$，每百万年的固体厚度。

\textbf{10.} 页岩的供应最快，白垩最慢——但三者之间
的差异仅约五分之一，远小于它们今日厚度所
暗示的程度。

\textbf{11.} \qty{800}{m}：其顶面今日位于地表，沉积
刚刚结束，因此其沉积时的厚度即为今日厚度。

\textbf{12.} 对页岩，由 $s(0, b) = 1306.7$ 按
\cref{met:b1:diagenesis-basins:decompaction}求解，得 $b = \qty{2124}{m}$：页岩
在 \qty{15}{\Ma} 时厚 \qty{2124}{m}。

\textbf{13.} 对白垩，由 $s(0, b) = 546.0$ 得
\qty{55}{\Ma} 时厚 \qty{1071}{m}。

\textbf{14.} 顶面位于 \qty{2124}{m} 处，由 $s(2124, b) = 546.0$ 得 $b =
\qty{2748}{m}$：白垩厚 \qty{624}{m}。

\textbf{15.} \qty{55}{\Ma} 时为 \qty{1071}{m}；$2124 + 624 = \qty{2748}{m}$，
在 \qty{15}{\Ma} 时。

\textbf{16.} 由今日厚度：砂岩 $800/15 = \qty{53.3}{m}$，
页岩 $1800/40 = \qty{45.0}{m}$，白垩 $600/20 = \qty{30.0}{m}$，每百万年
的速率。由沉积时的厚度：\qty{53.3}{m}，$2124/40 =
\qty{53.1}{m}$，$1071/20 = \qty{53.6}{m}$。第二组数据才是有意义的，
这一组——即海底增生的速率——显示出稳定的
速率，而今日厚度掩盖了这一点。

\textbf{17.} $1 - 1800/2124 = \qty{15}{\%}$.

\textbf{18.} 每个区间起止时的水深相同，因此所产生的
\omterm{def:b1:diagenesis-basins:accommodation}{可容纳空间}即基准面以上柱体的增长量：
自 75 至 \qty{55}{\Ma} 为 \qty{1071}{m}，$2748 - 1071 = \qty{1677}{m}$，
自 55 至 \qty{15}{\Ma}，$3200 - 2748 = \qty{452}{m}$，自 \qty{15}{\Ma} 以来。

\textbf{19.} $S = \Delta A - \Delta E$：$1677 - (-60) = \qty{1737}{m}$，自 55 至
\qty{15}{\Ma}；$452 - (-30) = \qty{482}{m}$，以来。

\textbf{20.} $1737/40 = \qty{43}{m}$，$482/15 = \qty{32}{m}$，每百万年
的速率。

\textbf{21.} $60/1737 = \qty{3.5}{\%}$，$30/482 = \qty{6.2}{\%}$。海平面在两个区间内均下降，
因此\omterm{def:b1:diagenesis-basins:eustasy}{全球海平面变化}阻碍了充填：基底须多\omterm{def:b1:diagenesis-basins:subsidence}{沉降}
这些量才能腾出同样的空间。\omterm{def:b1:diagenesis-basins:subsidence}{沉降}几乎提供了
全部\omterm{def:b1:diagenesis-basins:accommodation}{可容纳空间}。

\textbf{22.} 页岩与白垩在 \qty{15}{\Ma} 时合计厚 \qty{2748}{m}，
今日厚 \qty{2400}{m}：其\omterm{def:b1:diagenesis-basins:compaction}{压实}产生了 \qty{348}{m} 的
空间，占砂岩所充填 \qty{452}{m} 的 \qty{77}{\%}。

\textbf{23.} 每个岩层都因其之后的埋藏而减薄，减薄量
取决于其岩性及埋深；每个区间内产生的空间
必须用当时各岩层的厚度来度量，而
只有\omterm{def:b1:diagenesis-basins:compaction}{去压实}才能给出该厚度。即便如此，由此求得的\omterm{def:b1:diagenesis-basins:subsidence}{沉降}仍包括
沉积物自重引起的下沉，须扣除后才得到
\omterm{def:b1:diagenesis-basins:subsidence}{构造沉降}。

\textbf{24.} \qty{2124}{m}，约 \qty{2.1}{km}——比
井中今日所见的 \qty{1800}{m} 多 \qty{324}{m}。
\end{solution}
